% ============================================================
%  EXAMEN TEMA 3: DESPEJE DE FÓRMULAS Y APLICACIONES
%  Curso Introductorio — 3er Año
%  Duración: 90 minutos
%  Temas: Despeje lineal, despeje con exponentes y radicales,
%         geometría espacial (volúmenes) y densidad
%  Autor: Ing. Gabriel Astudillo
%  Nota: enunciado limpio (sin pistas ni desarrollo). Las
%  soluciones están en la "Clave de Respuestas" al final.
% ============================================================

\documentclass[12pt, letterpaper]{article}

% ── Paquetes ─────────────────────────────────────────────────

\usepackage{mathwizards-guia}
\mwguia{Examen — Tema 3: Despeje de Fórmulas y Aplicaciones}{Ing. Gabriel Astudillo $\cdot$ Curso 3er Año}

\begin{document}
% ============================================================

% ── Portada / Título ─────────────────────────────────────────
\mwportada{Examen del Tema 3}{Despeje de Fórmulas y Aplicaciones Integradoras}{Transposición, Geometría Espacial y Densidad}

\vspace{4pt}
\begin{cuadroinfo}
  \small
  \textbf{Nombre:} \rule{0.55\linewidth}{0.4pt} \quad \textbf{Fecha:} \rule{0.15\linewidth}{0.4pt} \\[6pt]
  \textbf{Tiempo disponible:} 90 minutos \quad$\bullet$\quad \textbf{Puntaje total:} 100 puntos \\[4pt]
  \textbf{Instrucciones:}
  \begin{enumerate}[leftmargin=*]
    \item Lee cada ejercicio con calma antes de resolverlo.
    \item Muestra \emph{todos} tus pasos; en los despejes, justifica cada operación inversa.
    \item Usa $\pi \approx 3{,}14$ cuando se requiera un valor numérico.
    \item Cuida las \textbf{unidades} y usa $1\,\text{mL} = 1\,\text{cm}^3$ y $1\,\text{L} = 1\,000\,\text{cm}^3$.
    \item No se permite el uso de calculadora.
    \item Escribe con letra clara y revisa tu examen antes de entregarlo.
  \end{enumerate}
\end{cuadroinfo}

\vspace{8pt}

% ============================================================
\section{Despeje de Fórmulas I: Variables Lineales (25 puntos)}
% ============================================================

\begin{enumerate}[leftmargin=*]
  \item \textbf{(10 pts)} Despeja la variable indicada:
        \[
          a)\ P = 2\ell + 2w \ \ (\text{despeja } w) \qquad
          b)\ C = \frac{5}{9}(F - 32) \ \ (\text{despeja } F)
        \]
        \[
          c)\ Q = m\,c\,\Delta T \ \ (\text{despeja } c) \qquad
          d)\ n = \frac{m}{M} \ \ (\text{despeja } M)
        \]

  \item \textbf{(8 pts)} Despeja:
        \[
          a)\ v_f = v_0 + a\,t \ \ (\text{despeja } v_0) \qquad
          b)\ \frac{1}{R_T} = \frac{1}{R_1} + \frac{1}{R_2} \ \ (\text{despeja } R_T)
        \]

  \item \textbf{(7 pts)} En la fórmula de la densidad $d = \dfrac{m}{V}$, despeja $V$. Luego, si $d = 8{,}96\;\text{g/cm}^3$ y $m = 448\;\text{g}$, calcula el volumen.
\end{enumerate}

\newpage

% ============================================================
\section{Despeje de Fórmulas II: Exponentes y Radicales (25 puntos)}
% ============================================================

\begin{enumerate}[leftmargin=*]
  \item \textbf{(12 pts)} Despeja la variable indicada:
        \[
          a)\ A = \pi r^2 \ \ (\text{despeja } r) \qquad
          b)\ v = \sqrt{2gh} \ \ (\text{despeja } h)
        \]
        \[
          c)\ E_c = \frac{1}{2}mv^2 \ \ (\text{despeja } v) \qquad
          d)\ c^2 = a^2 + b^2 \ \ (\text{despeja } a)
        \]

  \item \textbf{(7 pts)} Despeja $L$ de la fórmula del péndulo:
        \[
          T = 2\pi\sqrt{\frac{L}{g}}
        \]

  \item \textbf{(6 pts)} Despeja $x$ de $v_f^2 = v_0^2 + 2a\,x$. Luego calcula $x$ si $v_f = 20\;\text{m/s}$, $v_0 = 0$ y $a = 2\;\text{m/s}^2$.
\end{enumerate}

\newpage

% ============================================================
\section{Geometría Espacial Aplicada: Volúmenes (25 puntos)}
% ============================================================

\begin{enumerate}[leftmargin=*]
  \item \textbf{(8 pts)} Calcula el volumen de:
        \begin{enumerate}[label=\alph*)]
          \item una caja rectangular de $12\;\text{cm} \times 5\;\text{cm} \times 8\;\text{cm}$;
          \item una esfera de radio $3\;\text{cm}$;
          \item un cilindro de radio $2\;\text{cm}$ y altura $10\;\text{cm}$.
        \end{enumerate}

  \item \textbf{(8 pts)} Un cilindro tiene un volumen de $900\;\text{cm}^3$ y un radio de $5\;\text{cm}$. \textquestiondown Cuál es su altura?

  \item \textbf{(9 pts)} Un acuario tiene forma de paralelepípedo con dimensiones $80\;\text{cm} \times 40\;\text{cm} \times 50\;\text{cm}$.
        \begin{enumerate}[label=\alph*)]
          \item Calcula su volumen en $\text{cm}^3$ y en litros.
          \item Si se llena hasta $\dfrac{3}{4}$ de su capacidad, \textquestiondown cuántos litros de agua contiene?
        \end{enumerate}
\end{enumerate}

\newpage

% ============================================================
\section{La Densidad Absoluta (25 puntos)}
% ============================================================

\begin{enumerate}[leftmargin=*]
  \item \textbf{(8 pts)}
        \begin{enumerate}[label=\alph*)]
          \item Un bloque tiene una masa de $720\;\text{g}$ y un volumen de $900\;\text{cm}^3$. Calcula su densidad. \textquestiondown Flota en el agua?
          \item \textquestiondown Cuál es la masa de $300\;\text{cm}^3$ de aluminio, si su densidad es $2{,}7\;\text{g/cm}^3$?
        \end{enumerate}

  \item \textbf{(8 pts)} Una probeta contiene $60\;\text{mL}$ de agua. Al introducir una piedra, el nivel sube a $78\;\text{mL}$. Si la piedra tiene una masa de $45\;\text{g}$, calcula su densidad.

  \item \textbf{(9 pts)} Un objeto tiene forma de paralelepípedo con dimensiones $5\;\text{cm} \times 4\;\text{cm} \times 10\;\text{cm}$ y una masa de $1\,580\;\text{g}$.
        \begin{enumerate}[label=\alph*)]
          \item Calcula su densidad.
          \item Identifica el posible material (densidad del hierro: $7{,}87\;\text{g/cm}^3$).
          \item \textquestiondown Qué masa tendría un cubo de $1\;\text{m}$ de lado del mismo material? (Expresa el resultado en kilogramos.)
        \end{enumerate}
\end{enumerate}

\newpage

% ============================================================
\ifmwclaves
  \section{Clave de Respuestas}
  % ============================================================

  \subsection*{Sección 1: Despeje de Fórmulas I: Variables Lineales}

  \begin{enumerate}[leftmargin=*, label=\textbf{\arabic*.}]
    \item a) $P - 2\ell = 2w \Rightarrow w = \dfrac{P - 2\ell}{2}$. \quad
          b) $\dfrac{9C}{5} = F - 32 \Rightarrow F = \dfrac{9C}{5} + 32$. \quad
          c) $c = \dfrac{Q}{m\,\Delta T}$. \quad
          d) $M = \dfrac{m}{n}$.

    \item a) $v_0 = v_f - a\,t$. \quad
          b) $\dfrac{1}{R_T} = \dfrac{R_2 + R_1}{R_1 R_2} \Rightarrow R_T = \dfrac{R_1 R_2}{R_1 + R_2}$.

    \item $V = \dfrac{m}{d}$. \quad Sustituyendo: $V = \dfrac{448}{8{,}96} = 50\;\text{cm}^3$.
  \end{enumerate}

  \newpage

  \subsection*{Sección 2: Despeje de Fórmulas II: Exponentes y Radicales}

  \begin{enumerate}[leftmargin=*, label=\textbf{\arabic*.}]
    \item a) $r = \sqrt{\dfrac{A}{\pi}}$. \quad
          b) $v^2 = 2gh \Rightarrow h = \dfrac{v^2}{2g}$. \quad
          c) $v = \sqrt{\dfrac{2E_c}{m}}$. \quad
          d) $a = \sqrt{c^2 - b^2}$.

    \item Dividimos por $2\pi$: $\dfrac{T}{2\pi} = \sqrt{\dfrac{L}{g}}$. Elevamos al cuadrado: $\dfrac{T^2}{4\pi^2} = \dfrac{L}{g}$.
          Entonces $L = \dfrac{g\,T^2}{4\pi^2}$.

    \item $2a\,x = v_f^2 - v_0^2 \Rightarrow x = \dfrac{v_f^2 - v_0^2}{2a}$.
          Sustituyendo: $x = \dfrac{20^2 - 0^2}{2(2)} = \dfrac{400}{4} = 100\;\text{m}$.
  \end{enumerate}

  \newpage

  \subsection*{Sección 3: Geometría Espacial Aplicada: Volúmenes}

  \begin{enumerate}[leftmargin=*, label=\textbf{\arabic*.}]
    \item a) $V = 12 \times 5 \times 8 = 480\;\text{cm}^3$.
          \quad b) $V = \dfrac{4}{3}\pi(3)^3 = 36\pi \approx 113{,}1\;\text{cm}^3$.
          \quad c) $V = \pi(2)^2(10) = 40\pi \approx 125{,}7\;\text{cm}^3$.

    \item $V = \pi r^2 h \Rightarrow h = \dfrac{V}{\pi r^2} = \dfrac{900}{\pi(5)^2} = \dfrac{900}{25\pi} \approx 11{,}46\;\text{cm}$.

    \item a) $V = 80 \times 40 \times 50 = 160\,000\;\text{cm}^3 = 160\;\text{L}$.
          \quad b) $\dfrac{3}{4}(160) = 120\;\text{L}$.
  \end{enumerate}

  \newpage

  \subsection*{Sección 4: La Densidad Absoluta}

  \begin{enumerate}[leftmargin=*, label=\textbf{\arabic*.}]
    \item a) $d = \dfrac{720}{900} = 0{,}8\;\text{g/cm}^3 < 1\;\text{g/cm}^3 \Rightarrow$ \textbf{sí flota} en el agua.
          \quad b) $m = d \cdot V = 2{,}7 \times 300 = 810\;\text{g}$.

    \item $V = 78 - 60 = 18\;\text{mL} = 18\;\text{cm}^3$. \quad $d = \dfrac{45}{18} = 2{,}5\;\text{g/cm}^3$.

    \item a) $V = 5 \times 4 \times 10 = 200\;\text{cm}^3$; \quad $d = \dfrac{1\,580}{200} = 7{,}9\;\text{g/cm}^3$.
          \quad b) El valor $7{,}9\;\text{g/cm}^3$ es muy cercano al del \textbf{hierro} ($7{,}87\;\text{g/cm}^3$).
          \quad c) Un cubo de $1\;\text{m}$ de lado tiene $V = 10^6\;\text{cm}^3$, entonces $m = 7{,}9 \times 10^6\;\text{g} = 7\,900\;\text{kg}$.
  \end{enumerate}

\fi
\end{document}
